%% EXHALE: how molecular hydrogen is treated.
%% H2 is carried as ONE species with no (v,J) resolution, so every place that
%% needs a level distribution supplies one by assumption -- and the assumptions
%% differ from site to site.  This memo collects them, with the evidence each
%% rests on and the range over which it holds.
%% Author: Kwang-Il Seon
\documentclass[english,a4paper]{my_memo}
\synctex=1
\usepackage{booktabs}
\usepackage{array}
\usepackage{longtable}
\usepackage{babel}
\emergencystretch=3em

\title{Molecular hydrogen in EXHALE}
\author{Kwang-Il Seon}
\date{Last updated: \docmoddate}

\begin{document}
\maketitle

\begin{abstract}
EXHALE carries H$_2$ as a single species. The state vector holds one number
density for it, the network reacts it as one reagent, and no rovibrational
level is transported. Every quantity that depends on how the molecules are
distributed over $(v,J)$ therefore has that distribution supplied by
assumption --- and the assumption is not the same in every place. The infrared
line cooling is explicitly LTE and says so in its own function name; the
thermal dissociation and its three-body reverse rest on vibrational
equilibrium and are now built from one another through the equilibrium
constant, so that axis is consistent by construction; the Lyman--Werner
self-shielding factor imports the non-LTE distribution of a $10^{2}$~K
interstellar photodissociation front into a collision-dominated base above
$10^{3}$~K; and the H$_3^+$
cooling, alone among the four, carries an explicit departure factor. This memo
records each assumption, the published evidence behind it, the range over which
it holds, and where the layer EXHALE actually computes leaves that range.
\end{abstract}

\section{What is and is not carried}

H$_2$ enters the state vector as species \texttt{isp\_H2} of
\texttt{src/modules/init/species\_table.f90}, one of the four molecular
columns (\texttt{H2}, \texttt{H2p}, \texttt{H3p}, \texttt{HeHp}) that the
molecular option adds. It is written to \texttt{output/Ion\_species.txt}
and \texttt{output/Ion\_species\_adv.txt} as the number density of the
molecule; the four molecular columns are the last four of the file, but their
absolute position depends on which metal elements the run carries, so they
must be located from the \texttt{\# columns} schema header rather than
counted.

What the state vector does \emph{not} carry is any resolution of the
rovibrational ladder. There is one H$_2$ continuity equation, not one per
level, and no state-to-state collisional network. Nothing in the code
transports a vibrational temperature, an ortho/para ratio, or a level
population of any kind. The rovibrational ladder appears twice as
\emph{tabulated data} --- the partition function of the equilibrium constant
(\texttt{mol\_rates.f90}) and the line list of the infrared cooling
(\texttt{molecular\_infrared\_data.f90}) --- but in both places it is
evaluated at the local kinetic temperature and never solved for.

The consequence is the subject of this memo. A rate coefficient measured on
one distribution of levels does not transfer to another, and EXHALE uses
coefficients measured under at least three different distributions in the same
cell. Each of them is defensible where it is used, and the point of collecting
them is to say exactly where that is.

\subsection{Code map}

Table~\ref{tab:codemap} lists every file that touches H$_2$ and what it does
with it. Three groups are worth separating in advance, because the
level-distribution assumption enters each of them differently: the files that
supply a \emph{rate}, those that supply an \emph{emissivity}, and those that
supply a \emph{composition}.

\begin{table}[htbp]
\centering
\small
\begin{tabular}{@{}>{\raggedright\arraybackslash}p{0.29\textwidth}>{\raggedright\arraybackslash}p{0.62\textwidth}@{}}
\toprule
File & What it does with H$_2$ \\
\midrule
\multicolumn{2}{@{}l}{\textbf{\texttt{init/}}} \\
\texttt{species\_table.f90}
  & Declares \texttt{isp\_H2} and the nucleus weights ($2$ H nuclei per
    molecule) the element budget counts with \\
\texttt{set\_IC.f90}
  & Seeds the H$_2$ column from the chemical-equilibrium fit at the local
    pressure and $T_0$, capped by the available neutral H \\
\addlinespace
\multicolumn{2}{@{}l}{\textbf{\texttt{lower\_atmosphere/}}} \\
\texttt{mol\_rates.f90}
  & The Koskinen et al.\ (2022) Table-1 network, plus the H$_2$
    thermochemistry: partition function, equilibrium constant, three-body
    recombination \\
\texttt{lower\_column.f90}
  & \texttt{q\_h2\_equilibrium}, the chemical-equilibrium mixing ratio, and
    the mean molecular weight of the analytic column built from it \\
\texttt{lyman\_werner.f90}
  & Solomon-process photodissociation and the Draine \& Bertoldi (1996)
    self-shielding function \\
\texttt{molecular\_infrared\_\allowbreak cooling.f90},
\texttt{\dots\_data.f90}
  & The H$_2$ quadrupole and magnetic-dipole line cooling, and the H$_2$O and
    CO bands beside it \\
\texttt{h3p\_cooling.f90}
  & H$_3^+$ infrared cooling; H$_2$ enters as the collider density of the
    non-LTE departure factor \\
\addlinespace
\multicolumn{2}{@{}l}{\textbf{\texttt{nonlinear\_system\_solver/}}} \\
\texttt{System\_HeH\_mol.f90},
\texttt{System\_HeH\_mol\_\allowbreak metals.f90}
  & The coupled ionization system: unknown $x(4) = 2n_{\rm H_2}/n_{\rm H}$,
    the H nuclei bound in H$_2$ \\
\texttt{constrained\_chemical\_\allowbreak equilibrium.f90}
  & The same network in logarithmic densities, with a continuation in the
    radiation field seeded from the H$_2$ dissociation equilibrium \\
\addlinespace
\multicolumn{2}{@{}l}{\textbf{\texttt{functions/}}} \\
\texttt{composition.f90}
  & \texttt{h2\_mixing\_ratio\_base}, \texttt{h2\_bound\_fraction},
    \texttt{comp\_ntot\_bc}: the H nuclei bound into H$_2$ leave the base
    particle count \\
\texttt{cross\_sec.f90}
  & \texttt{sigma\_H2}, the photoionization cross section above $15.4$~eV \\
\texttt{binary\_element\_\allowbreak diffusion.f90}
  & H$_2$ as a diffusion carrier of the hydrogen component, with mass $2$ and
    its own polarizability \\
\addlinespace
\multicolumn{2}{@{}l}{\textbf{\texttt{radiation/} and \texttt{post\_process/}}} \\
\texttt{util\_ion\_eq.f90}
  & \texttt{fuv\_lw\_photon\_field}: the shared $912$--$1110$~\AA{} beam, the
    H$_2$ column, and the band fraction the pumping lines remove \\
\texttt{Cool\_coeff.f90}
  & \texttt{ioniz\_HeI23S\_H2}, the He($2^3S$) $+$ H$_2$ Penning and
    associative ionization rate \\
\texttt{post\_process\_adv.f90}
  & Does \emph{not} carry H$_2$: the \texttt{\_adv} reconstruction is
    molecule-free by construction, and says so in its header \\
\bottomrule
\end{tabular}
\caption{Where H$_2$ appears. Paths are relative to \texttt{src/modules/}.}
\label{tab:codemap}
\end{table}

\subsection{The layer the assumptions are evaluated in, and what the numbers
below are read from}
\label{sec:layer}

Every claim about whether a fit is used inside or outside its range needs a
statement of the conditions EXHALE actually reaches. The reference case here
is \texttt{mol\_lyman\_werner} of the regression matrix --- the hot-Uranus
molecular base with \texttt{Stellar LW flux [erg/cm2/s]: 343.0} --- because
it is the only default case that carries the Lyman--Werner channel.
Measured on its \texttt{output/Lyman\_Werner.txt} and
\texttt{output/Ion\_species.txt}:

\begin{center}
\small
\begin{tabular}{@{}lllll@{}}
\toprule
Layer definition & cells & $r/R_{\rm p}$ & $T$ [K] & $n_{\rm H}$ [cm$^{-3}$] \\
\midrule
$x_{\rm H_2} > 0.5$   & 227 & $0.9998$--$1.127$ & $1149$--$1307$ & $1.67\times10^{12}$--$9.27\times10^{13}$ \\
$x_{\rm H_2} > 0.01$  & 251 & $0.9998$--$1.172$ & $945$--$1307$  & $6.24\times10^{11}$--$9.27\times10^{13}$ \\
$x_{\rm H_2} > 10^{-6}$ & --- & $0.9998$--$3.853$ & $945$--$2430$ & $1.98\times10^{7}$--$9.27\times10^{13}$ \\
\bottomrule
\end{tabular}
\end{center}

\noindent
with $x_{\rm H_2} = 2n_{\rm H_2}/n_{\rm H}$; the H$_2$ density itself peaks at
$4.63\times10^{13}$~cm$^{-3}$ in the base cell and falls to
$4.42\times10^{9}$~cm$^{-3}$ at the $1\%$ edge. The three rows matter because they
give three different answers. The often-quoted range $945$--$2430$~K is the
extent of gas that carries \emph{any} H$_2$ at all, and its upper half is
gas where H$_2$ is a $10^{-5}$ trace; the layer where H$_2$ is actually a
constituent spans $945$--$1307$~K. Statements below use the narrow definition
unless they say otherwise.

\paragraph{These are relaxation snapshots, not converged solutions.}
Six of the eight default regression cases --- every molecular one, plus
\texttt{lower\_profile} --- are pinned at $12{,}000$ steps by a
\texttt{maxsteps} file and stop there rather than at a convergence criterion.
\texttt{mol\_lyman\_werner} ends at \texttt{du} $= 2.24$, i.e.\ a $224\%$
radial spread in the mass flux $\rho v r^2$, three orders of magnitude from
the $10^{-3}$ threshold a converged run meets; only the two \texttt{wasp}
cases converge. Every profile number in this memo therefore describes a
relaxation snapshot. That is adequate for the purpose --- the question is
which decade of density and temperature a fit is being asked about, not the
fourth digit --- but it is not adequate for quoting a mass-loss rate, and none
is quoted here.

\paragraph{And the snapshots do not currently match the goldens.} At the time
of writing, all eight default cases differ from
\texttt{backup/regression/golden/} in \texttt{Hydro\_ioniz.txt} and
\texttt{Ion\_species.txt}. The working tree carries uncommitted source
changes, including the one described in \S\ref{sec:diss}, and the goldens
predate them. The numbers above are from the case output directories as they
stand, not from the goldens.

\section{The level-distribution assumptions}

\subsection{Infrared line cooling: LTE, stated explicitly}
\label{sec:ir}

\texttt{molecular\_infrared\_cooling.f90} evaluates the H$_2$ line emission
as a Boltzmann sum,
%
\begin{equation}
  E_{\rm H_2}(T) = \sum_l f_u(T)\, A_l\, \Delta E_l ,
  \qquad
  f_u(T) = \frac{g_u\,e^{-T_u/T}}{Z(T)},
  \qquad
  Z(T) = \sum_{\rm lev} g\, e^{-T_{\rm lev}/T},
\end{equation}
%
and names the routine that returns it \texttt{h2\_line\_emission\_lte}. There
is no ambiguity about the assumption here; the question is only where it
holds.

The line list is Roueff et al.\ (2019), A\&A 630, A58, table 2 (electric
quadrupole plus magnetic dipole), and the level ladder comes from the same
table. As carried in \texttt{molecular\_infrared\_data.f90} it is
\texttt{n\_h2\_lev} $= 302$ levels and \texttt{n\_h2\_line} $= 1833$ lines.

Rotational levels are not the constraint: they thermalize at densities far
below anything in this layer. The vibrational levels are. The code records the
comparison at its own site: $n_{\rm crit}({\rm H_2}\,v=1) \sim
10^{9}$--$10^{11}$~cm$^{-3}$ against $10^{13}$--$10^{14}$~cm$^{-3}$ at the
base --- and the measured base density of the molecular gate case,
$n({\rm H_2}) = 4.63\times10^{13}$~cm$^{-3}$, sits in that range.

Above the base the density falls, but so does the temperature, and the
vibrational bands then carry a small share of the emission. That is the
quantitative statement that makes the LTE overestimate harmless where LTE is
least justified, and it can be evaluated directly from the line list: sum the
emission of the lines whose upper level has $v \geq 1$ and divide by the
total. Table~\ref{tab:vibshare} does so.

\begin{table}[h]
\centering
\begin{tabular}{@{}lll@{}}
\toprule
$T$ [K] & $E_{\rm H_2}(T)$ [erg s$^{-1}$ molecule$^{-1}$]
        & share from $v_u \geq 1$ \\
\midrule
$400$  & $3.034\times10^{-23}$ & $0.72\%$ \\
$500$  & $7.717\times10^{-23}$ & $5.70\%$ \\
$600$  & $1.789\times10^{-22}$ & $18.2\%$ \\
$1000$ & $2.726\times10^{-21}$ & $65.7\%$ \\
\bottomrule
\end{tabular}
\caption{Share of the LTE H$_2$ line emission carried by transitions out of a
vibrationally excited upper level, evaluated on the implemented line list.
Outside the tabulated range the share is $0.016\%$ at $300$~K and $81.5\%$ at
$2000$~K. The level table carries no $(v,J)$ index, so the vibrational
assignment was recovered from the Roueff et al.\ (2019) source table and
matched back onto the $1833$ selected lines; classifying instead by
$v_u \neq v_l$ (a band transition rather than a pure rotational one) changes
the shares by less than $0.1$ percentage point up to $1000$~K.}
\label{tab:vibshare}
\end{table}

The reading is that the LTE assumption is safe at the base, where the density
supports it, and that where the density no longer supports it the vibrational
bands --- the levels the density fails to thermalize --- have already stopped
carrying the emission. Below about $500$~K the error LTE can make is bounded
by a few per cent of the total, because that is all the vibrational levels
contribute.

\subsection{Thermal dissociation and three-body recombination: vibrational
equilibrium}
\label{sec:diss}

The pair R12 (H$_2$ + M $\rightarrow$ H + H + M) and R15 (H + H + M
$\rightarrow$ H$_2$ + M) is one channel run in two directions. Since the
change recorded below, \texttt{mol\_rates.f90} builds R12 from R15 through the
equilibrium constant,
%
\begin{equation}
  k_{\rm diss}(T) = \frac{k_{\rm rec}(T)}{K_{\rm eq}(T)},
  \qquad
  K_{\rm eq}(T) = \frac{n_{\rm H_2}}{n_{\rm H}^2}\ [\mathrm{cm^3}],
\end{equation}
%
rather than transcribing an independent dissociation fit.

That construction is what makes this axis consistent. $K_{\rm eq}$ is built in
\texttt{h2\_thermochemistry\_init} from the translational partition functions
and the \emph{rovibrational} internal partition function \texttt{q\_rovib\_H2},
summed over the Morse ladder of H$_2$ $X\,^1\Sigma_g^+$ with the ortho/para
nuclear-spin weights. It therefore contains vibrational excitation explicitly:
both sides of the equilibrium refer to a vibrationally equilibrated gas, and
the reverse rate built from it inherits that and nothing else.

The measured coefficient it is anchored on refers to the same state, and the
source says so. Breshears \& Bird (1973), 14th Symposium (International) on
Combustion, 211, read their initial rate off the minimum of the postshock
density-gradient dip, which they identify as ``the rate corresponding to
conditions of vibrational equilibrium in the thermal sense but with negligible
extent of dissociation''. The same paper bounds where its rate law applies:
the ``observation of induction periods for diatomic dissociation clearly
indicates that the assumptions implicit in Eqs.\ (2)--(4) are not valid during
the very early course of the reaction; nevertheless, they may be applicable
during the `steady state' which follows the induction period.'' A
hydrodynamic model that resolves nothing faster than a vibrational relaxation
time is in that steady state throughout. Their Eq.\ (3), $k_d/k_r = K_{\rm
eq}$, states the detailed-balance relation outright, so deriving one direction
from the other is the relation the measured coefficient already carries rather
than an assumption added here.

The critical density for vibrational thermalization, $10^{4}$--$10^{8}$~
cm$^{-3}$ depending on the collider, sits many orders below the
$10^{13}$--$10^{15}$~cm$^{-3}$ of the base, which is the independent check
that the premise holds where the rate is evaluated.

\paragraph{Status of this change.} R12 as detailed balance is present in the
working tree and \emph{not} in the last commit: \texttt{git show
HEAD:src/modules/lower\_atmosphere/mol\_rates.f90} still carries the
independent Baulch et al.\ (1992) fit $1.5\times10^{-9}e^{-48350/T}$, and
\texttt{docs/Update\_EXHALE.md} ends at section 114 with no entry for the
change. The changelog is therefore behind the code on this point.

\subsection{Lyman--Werner self-shielding: the weakest of the four}
\label{sec:lw}

\texttt{lyman\_werner.f90} takes the self-shielding factor from Draine \&
Bertoldi (1996), ApJ 468, 269 (hereafter DB96), their eq.~(37),
%
\begin{equation}
  f_{\rm shield}(N_{\rm H_2}) =
    \frac{0.965}{(1 + x/b_5)^2}
  + \frac{0.035}{(1+x)^{1/2}}\,
    \exp\!\left[-8.5\times10^{-4}(1+x)^{1/2}\right],
  \qquad x = \frac{N_{\rm H_2}}{5\times10^{14}\ \mathrm{cm^{-2}}},
\end{equation}
%
with $b_5 = b/10^5$~cm~s$^{-1}$ and $b$ the thermal Doppler parameter
$(2kT/m_{\rm H_2})^{1/2}$.

This is where the imported level distribution is furthest from ours. DB96
built the fit for interstellar photodissociation regions: their stationary
fronts run at $n_{\rm H} = 10^{2}$~cm$^{-3}$ with the temperature parameter
$T_0 = 200$~K (their Figs.~3--6 and 8--11; Figs.~1 and 2 are $100$ and
$300$~K LTE comparisons, and Fig.~7, the one this section leans on, is at
$T = 10^{2}$~K with no dust). In that gas the level distribution of the
\emph{shielded} H$_2$ is not thermal at all --- it is set in non-LTE by
ultraviolet pumping and by formation on grains. Neither $T$ nor $n_{\rm H}$ is
an argument of the fit ($T$ enters only through $b$), so what those conditions
really fix is the rovibrational population of the absorbing H$_2$, which is
the one ingredient of DB96's calculation that cannot be rescaled from outside.
Normalizing $f_{\rm shield}$ to the $N = 0$ rate carries with it the
assumption that this distribution does not change with depth.

EXHALE evaluates it in a collision-dominated planetary base at
$945$--$1307$~K and $n_{\rm H}$ up to $9.3\times10^{13}$~cm$^{-3}$, eleven
orders of magnitude denser than the fronts the fit was built on and five to
six times hotter.

\paragraph{What partly resolves it, and what the paper actually says.} DB96's
Fig.~7 compares the self-shielding produced by three different level
distributions --- a $100$~K LTE population and two ultraviolet-pumped non-LTE
ones at $\chi = 1$ and $\chi = 10$. Its \emph{caption} states that ``the three
self-shielding factors are essentially identical for $N({\rm H_2}) \gtrsim
2\times10^{17}$~cm$^{-2}$''.

The \emph{body} text of the same section is more careful, and it is the
better ground to stand on:

\begin{quote}\small
``We see that for $10^{14} \lesssim N_2 \lesssim 10^{18}$~cm$^{-2}$,
$f_{\rm shield}$ does show a dependence on the details of the H$_2$
rovibrational distribution. When there is higher excitation, the pumping
occurs via a larger number of lines, so that the H$_2$ remains optically thin
a bit longer, and self-shielding is less effective until one enters the
heavily damped regime at $N_2 \gtrsim 10^{18}$~cm$^{-2}$. Overall, however,
the self-shielding function $f_{\rm shield}(N_2)$ is relatively insensitive to
the effects of UV pumping with $\chi \lesssim 10^{3}$.''
\end{quote}

\noindent
So the authors do \emph{not} claim distribution-independence at $10^{17}$; they
claim a dependence up to $10^{18}$ and its disappearance in the heavily damped
regime beyond. Our molecular layer sits at $10^{21}$~cm$^{-2}$ and above,
three decades inside that regime, which is a stronger position than the
caption alone would give.

\paragraph{What is left.} The insensitivity was demonstrated among \emph{cold}
distributions: a $100$~K LTE population against UV-pumped ones, all at
$T = 10^{2}$~K and $n_{\rm H} = 10^{2}$~cm$^{-3}$. None of the three is a
$10^{3}$~K collisionally equilibrated population, and no distribution at our
temperature was tested. The physical argument for the damped regime --- that
once the pumping lines are on their damping wings the equivalent width stops
caring which lines carry the absorption --- is what makes the transfer
plausible, not a demonstration. That is the honest statement of the position.

\paragraph{Two further axes.} The \emph{column} is where the layer leaves what
was checked --- though it is worth being precise about what that means, since
DB96 put no upper bound on eq.~(37) at all. Their \S5.2 says only that it ``does an
excellent job in reproducing the initial rapid decline in $f_{\rm shield}$ for
$10^{14} < N_2 < 10^{17}$~cm$^{-2}$, the `square root' behavior for $10^{17} <
N_2 < 10^{20}$~cm$^{-2}$, and the rapid falloff due to line overlap for
$N_2 > 10^{20}$~cm$^{-2}$''. What \emph{is} bounded is the comparison against
exact multiline transfer, which in the same section reaches ``the largest
column densities considered, $N_2 = 3\times10^{21}$~cm$^{-2}$, where line
overlap effects suppress the pumping rates by a factor $\sim\!10$''.

The base cell of \texttt{mol\_lyman\_werner} settles at $N_{\rm H_2} =
4.361\times10^{21}$~cm$^{-2}$ and passes through $5.31\times10^{21}$ while
relaxing, which the run reports once at the crossing. So the deepest cells
use the fit $45\%$ past the largest column at which anyone checked it ---
undemonstrated rather than outside a stated range, which is a weaker
indictment but still a real one, since the regime beyond $10^{20}$ is exactly
the line-overlap regime where the fit is doing its hardest work. There
$f_{\rm shield}$ is $9.5\times10^{-7}$ and falling as $N^{-3/4}$.

The \emph{Doppler parameter} enters only through the saturated core, the
$(1+x/b_5)^{-2}$ term. Our thermal $b$ is $2.79$~km~s$^{-1}$ at $945$~K and
$4.48$~km~s$^{-1}$ at $2430$~K, inside the $1$--$10$~km~s$^{-1}$ span the form
was built to cover. But at the columns of this layer the saturated term has
collapsed --- at $N_{\rm H_2} = 4.4\times10^{21}$ it contributes
$\sim\!10^{-10}$ against the $9.5\times10^{-7}$ of the damping-wing term ---
so $f_{\rm shield}$ is the same to seven digits at both $b$ values. The $b$
axis is neither tested nor, here, used: what the layer's shielding actually
rests on is the second term of eq.~(37), the one with no $b$ in it.

\subsection{H$_3^+$ cooling: a departure factor, not LTE}
\label{sec:h3p}

The one place where EXHALE does not assume LTE is the H$_3^+$ infrared
cooling, and it is worth stating because it runs the other way from
\S\ref{sec:ir}. \texttt{h3p\_cooling.f90} multiplies the Miller et al.\
(2013) LTE emission fit by the departure factor $s(T, n_{\rm H_2})$ of their
Table~6, bilinearly interpolated in $(T, \log_{10} n_{\rm H_2})$, clamped at
the table edges and set to $1$ for $n_{\rm H_2} \geq 10^{14}$~cm$^{-3}$. Here
H$_2$ enters as the collider density that sets how far the H$_3^+$ levels
depart from a Boltzmann distribution.

The net exchange is written in Kirchhoff form,
%
\begin{equation}
  \Lambda_{\rm net} = n_{\rm H_3^+}\, 4\pi\, s(T, n_{\rm H_2})
     \left[ E(T) - W\,E(T_{\rm rad}) \right],
\end{equation}
%
with $W$ the dilution of the field the unresolved lower atmosphere presents,
switched on by the input key \texttt{Base IR field} (default \texttt{False}).
The departure factor multiplies both terms, which it must: $s<1$ applied to
the emission alone would heat a gas hotter than the field, and the form above
is exactly zero at $T = T_{\rm rad}$, $W = 1$.

The factor is doing real work rather than sitting at its LTE value: evaluated
over the $x_{\rm H_2} > 0.01$ layer of \texttt{mol\_lyman\_werner}, $s$ runs
from $0.996$ in the base cell ($n_{\rm H_2} = 4.63\times10^{13}$~cm$^{-3}$)
down to $0.438$ at the top of that layer ($4.42\times10^{9}$~cm$^{-3}$,
$r = 1.172\,R_{\rm p}$). Assuming LTE there would overstate the H$_3^+$
cooling by a factor $2.3$.

The table is a vacuum departure factor --- collisions against spontaneous
decay, no incident field --- so a field strong enough to matter pumps the
levels back toward LTE and the true departure is smaller than $s$. The
treatment therefore underestimates the magnitude of the exchange, not its
sign or its fixed point.

For comparison, Frelikh \& Murray-Clay (2026), ApJ 996, 96 treat the same
coolant in pure LTE. On this one channel EXHALE carries the more complete
physics.

\subsection{The chemical-equilibrium mixing ratio}
\label{sec:qh2}

\texttt{q\_h2\_equilibrium} (\texttt{lower\_column.f90}) is the
Koskinen et al.\ (2022) Eq.~(11) fit to the Visscher et al.\ (2006)
chemical-equilibrium H$_2$/H partition,
%
\begin{equation}
  q_{\rm H_2} = \frac{1.9845 + 10^{u} - \sqrt{10^{u}(3.9690 + 10^{u})}}{2.3670},
  \qquad
  u = -\frac{23672}{T} - \log_{10} p[\mathrm{bar}] + 6.2645 ,
\end{equation}
%
with limits $1.9845/2.3670 = 0.83840$ as $u \rightarrow -\infty$ (fully
molecular, the fit's approximation of the solar-composition volume mixing
ratio $0.5/(0.5+{\rm He/H}) = 0.863$) and $0$ as $u \rightarrow +\infty$.

It is used in four places: the mean molecular weight of the analytic lower
column, the base H$_2$ mixing ratio when no \texttt{base.inp} supplies a
photochemical one, the initial-condition seed of the H$_2$ column, and the
$\lambda \rightarrow 0$ end of the continuation path in
\texttt{constrained\_chemical\_equilibrium.f90}.

The assumption here is not about level populations but about the chemistry:
this is \emph{thermochemical} equilibrium, and it knows nothing about
photodissociation. It is therefore not a statement about an irradiated wind,
and the code says so at its own site --- in chemical equilibrium
$q_{\rm H_2}(1\,\mu\mathrm{bar})$ stays large up to $T \sim 2000$~K and the
gas is fully atomic only above $\sim 2400$~K, while photochemistry plus the
hotter actual base temperature dissociate H$_2$ well below that. Using it as a
\emph{seed} for a solve that then applies the full radiation field is exactly
the use it supports; using it as a discriminant for whether a cell should be
molecular is not. Section 114 of \texttt{docs/Update\_EXHALE.md} records an
attempt in that direction being withdrawn, together with the repair of an
asymptotic guard in the same function that had been inverted since
transcription and handed every cell below about $529$~K a fully atomic state
labelled ``the dense molecular limit''.

\subsection{Three sites that inherit a ground-state assumption silently}

Three further quantities are measured on ground-state H$_2$ and used without
a level-distribution correction. None is flagged in the code as such, and
none is likely to matter at the level the three above do, but the list should
be complete.

\begin{itemize}
\item \textbf{Photoionization.} \texttt{sigma\_H2} in
  \texttt{functions/cross\_sec.f90} is the Yan, Sadeghpour \& Dalgarno (1998),
  ApJ 496, 1044 fit (their Eqs.\ 17--19), threshold $15.4$~eV, with the
  near-threshold resonance structure smoothed as in the source fit. It is a
  cross section for H$_2$ in its ground vibrational state; a vibrationally hot
  molecule has a lower threshold and different near-edge structure.
\item \textbf{Diffusion.} \texttt{binary\_element\_diffusion.f90} gives H$_2$
  the static dipole polarizability $\alpha = 5.315\,a_0^3$
  (Schwerdtfeger \& Nagle 2019) for the ion-neutral coefficient. This is the
  $v=0$ value; polarizability rises with vibrational excitation.
\item \textbf{Metastable-helium Penning ionization.}
  \texttt{ioniz\_HeI23S\_H2} in \texttt{radiation/Cool\_coeff.f90} is the
  Maxwell--Boltzmann average of the Cohen \& Lane (1977) cross sections, taken
  from the Garcia Munoz (2025) network file, valid over $500$--$10^4$~K
  following its tabulation. It is a thermal average over relative translational
  velocity with the internal state of the H$_2$ fixed.
\end{itemize}

\section{The third body}

R12, R13 and R15 are three-body reactions written as a two-body coefficient
times a density. Koskinen et al.\ (2022) print that density as ``$n$'' and do
not say which particles it counts; their models are H$_2$/H-dominated Jovian
and Neptunian envelopes, where M is H$_2$ and H. EXHALE passes the total
gas-particle density with electrons excluded
(\texttt{ion\_cell\_state\%ntot}), so every heavy particle is treated as an
equally efficient third body.

The coefficient this assumes uniform is the M~$=$~H$_2$ one: R15 carries
Cohen \& Westberg's (1983) recommended $k_1({\rm H_2}) = 2.8\times10^{-31}
T^{-0.6}$~cm$^6$~s$^{-1}$ over $50$--$5000$~K, labelled M~$=$~H$_2$ on their
data sheet and separate from their $k_1({\rm H})$.

Breshears \& Bird measured all three of their colliders in one experiment with
one analysis, so their dissociation coefficients are the clean statement of
how much the third body matters:
%
\begin{align}
  k_d(\mathrm{M}={\rm Ar}) &= 1.55\times10^{-10}\,e^{-44736/T}, \nonumber\\
  k_d(\mathrm{M}={\rm H_2}) &= 5.48\times10^{-9}\,e^{-52989/T},
  \qquad [\mathrm{cm^3\,s^{-1}}] \\
  k_d(\mathrm{M}={\rm H})  &= 3.52\times10^{-9}\,e^{-43881/T}. \nonumber
\end{align}
%
Evaluated here, $k_d({\rm H})/k_d({\rm H_2})$ is $8.67$ at $3500$~K, $6.26$ at
$4000$~K, $3.97$ at $5000$~K and $2.01$ at $8000$~K --- a factor of a few
across their measured $3500$--$8000$~K range. Cohen \& Westberg's independent
\emph{recombination} pair, $k_1({\rm H}) = 8.8\times10^{-33}$ against
$k_1({\rm H_2}) = 2.8\times10^{-31}T^{-0.6}$, gives $4.21$ and $5.21$ at
$3500$ and $5000$~K, which is the agreement detailed balance requires of two
independent sources inside their ranges.

Below the measured range the two Arrhenius forms diverge, because their
activation temperatures differ by $9108$~K: the ratio reaches $13.4$ at
$3000$~K, $61.0$ at $2000$~K and $5.80\times10^{3}$ at $1000$~K, against
$1.98$ from Cohen \& Westberg at $1000$~K. Take a factor of $2$ to $5$ as the
collider spread where H$_2$ exists, and read the $5.8\times10^{3}$ as the
cost of extrapolating an Arrhenius fit rather than as physics.

\paragraph{Helium is the collider nobody measured.} The equal-efficiency
assumption is weakest in two places, and only one of them is quantified above.
Above the H$_2$ front atomic H dominates the collider population, and the
factor of $2$--$5$ is the size of the error there --- carried, at least, by a
vanishing H$_2$ density. In a helium-rich envelope M is helium, and there is
no measurement to compare against at all. Breshears \& Bird shocked H$_2$ in
argon and in xenon ($5.0$, $10.0$, $20.0\%$ H$_2$ in Ar; $5.0$, $10.0$,
$20.2\%$ in Xe) and extracted M $=$ Ar, Xe, H$_2$ and H; helium appears in
their paper only as a shock-driver gas and in ``helium--neon laser''. None of
the sources behind R12, R13 or R15 supplies a helium efficiency, and none is
applied. On general grounds a monatomic third body has no internal modes to
take up the released energy and should be \emph{less} efficient than H$_2$, so
in a helium-dominated base these three rates are upper bounds;
\texttt{docs/molecular\_chemistry\_audit\_he\_rich.md} \S4.1 puts the size of
a plausible factor at order $0.5$ and records it as a guess rather than a
measurement. Nothing in this repository supports a helium efficiency near
unity, and no such number should be quoted from it.

\section{Validity ranges}
\label{sec:validity}

Table~\ref{tab:validity} collects every range in one place: what each
treatment stands on, and what EXHALE asks of it. The evaluated column is the
$x_{\rm H_2} > 0.01$ layer of \S\ref{sec:layer} unless stated, with the
caveats recorded there about relaxation snapshots.

\begin{longtable}{@{}>{\raggedright\arraybackslash}p{0.185\textwidth}>{\raggedright\arraybackslash}p{0.255\textwidth}>{\raggedright\arraybackslash}p{0.195\textwidth}>{\raggedright\arraybackslash}p{0.245\textwidth}@{}}
\caption{What each treatment stands on and what it is asked to do.}
\label{tab:validity}\\
\toprule
Treatment & Range of the data & Range evaluated & Verdict \\
\midrule
\endfirsthead
\multicolumn{4}{@{}l}{\emph{\tablename~\thetable\ (continued)}}\\
\toprule
Treatment & Range of the data & Range evaluated & Verdict \\
\midrule
\endhead
\bottomrule
\endlastfoot
H$_2$ LTE line emission
  & Vibrational thermalization needs
    $n \gg n_{\rm crit}(v{=}1) \sim 10^{9}$--$10^{11}$~cm$^{-3}$
  & $n({\rm H_2}) = 4.4\times10^{9}$--$4.6\times10^{13}$~cm$^{-3}$
  & Inside at the base; aloft the vibrational bands carry $<6\%$ of the
    emission below $500$~K, which bounds the error \\
\addlinespace
H$_2$O and CO band cross sections
  & Tabulated $50$--$2000$~K, clamped outside
  & $945$--$1307$~K in the molecular layer
  & Inside. Both molecules are largely dissociated above $2000$~K in any case
    \\
\addlinespace
R15 three-body recombination, $k_1({\rm H_2}) = 2.8\times10^{-31}T^{-0.6}$
  & Cohen \& Westberg (1983), recommended $50$--$5000$~K; $\pm0.2$ in
    $\log k$ at $300$~K rising to $\pm0.4$ at $5000$~K
  & $945$--$1307$~K
  & Inside \\
\addlinespace
R12 thermal dissociation, $k_{\rm rec}/K_{\rm eq}$
  & Inherits R15's range; $K_{\rm eq}$ tabulated $100$--$20000$~K and clamped
  & $945$--$1307$~K
  & Inside. Against the Baulch et al.\ (1992) fit the ratio is $0.97$--$1.18$
    over that fit's own $2500$--$8000$~K and departs only below it \\
\addlinespace
Third body M $=$ everything
  & Measured for M $=$ H$_2$, H, Ar, Xe (Breshears \& Bird 1973);
    no helium coefficient exists
  & Total heavy-particle density
  & H$_2$-dominated base: inside. Above the front, M is atomic H, a factor
    $2$--$5$. Helium: unquantified \\
\addlinespace
Lyman--Werner $f_{\rm shield}$, column axis
  & No upper bound stated. Checked against exact transfer to
    $3\times10^{21}$~cm$^{-2}$
  & up to $4.36\times10^{21}$, and $5.31\times10^{21}$ in transient
  & \textbf{Past what was checked}, by $45\%$ in column, in the line-overlap
    regime \\
\addlinespace
Lyman--Werner $f_{\rm shield}$, Doppler axis
  & Every exact multiline check DB96 ran was at $b = 3$~km~s$^{-1}$; the
    $1$--$10$~km~s$^{-1}$ axis is constructed
  & thermal $b = 2.79$--$4.48$~km~s$^{-1}$
  & Bracketing their one tested value, but the axis is untested. At our
    columns the $b$-dependent term is negligible anyway \\
\addlinespace
Lyman--Werner $f_{\rm shield}$, level distribution
  & $T = 10^{2}$~K, $n_{\rm H} = 10^{2}$~cm$^{-3}$, non-LTE (UV pumping and
    grain formation); Fig.~7's three distributions
  & $945$--$1307$~K, $n_{\rm H}$ up to $9.3\times10^{13}$~cm$^{-3}$,
    collisional
  & \textbf{Outside.} DB96 put the distribution dependence at
    $N_2 \lesssim 10^{18}$ and our columns are three decades above it, but no
    $10^{3}$~K distribution was tested \\
\addlinespace
Band fraction $A$ from DB96 eq.~(39)
  & Their own ceiling, \S4.3 eq.~(29), is
    $W_{\max} \approx \ln(1110/912) \approx 0.2$, i.e.\ $A = 1$
  & asymptote $1.0843$ at $b = 2.79$, $1.1050$ at $4.48$~km~s$^{-1}$
  & \textbf{Above the authors' own cap} by $8.4$--$10.5\%$; clamped at $1$ in
    the code \\
\addlinespace
H$_3^+$ emission fit $E(T)$
  & Miller et al.\ (2013) Table 5, piecewise $30$--$5000$~K; frozen above
  & $945$--$1307$~K
  & Inside. The published segments do not join; each is used on its own range
    \\
\addlinespace
H$_3^+$ departure factor $s(T, n_{\rm H_2})$
  & Table 6 grid $T = 300$--$5000$~K, $\log_{10} n_{\rm H_2} = 6$--$14$;
    clamped, $s = 1$ above $10^{14}$
  & $T = 945$--$1307$~K, $n_{\rm H_2} = 4.4\times10^{9}$--$4.6\times10^{13}$
  & Inside; $s = 0.44$--$0.996$ across the layer \\
\addlinespace
\texttt{q\_h2\_equilibrium}
  & Koskinen et al.\ (2022) Eq.~(11), validated against NASA CEA at
    $1000$--$2500$~K; thermochemical equilibrium only
  & Seed and base-composition use only
  & Inside as a seed. Not applicable to an irradiated wind: it has no
    photodissociation \\
\addlinespace
\texttt{sigma\_H2} photoionization
  & Yan et al.\ (1998) Eqs.~17--19, $E \geq 15.4$~eV, ground vibrational
    state, resonances smoothed
  & the full stellar spectrum above threshold
  & Inside in energy; the vibrational-state assumption is unquantified \\
\addlinespace
He($2^3S$) $+$ H$_2$ ionization
  & Cohen \& Lane (1977) via Garcia Munoz (2025), tabulated
    $500$--$10^4$~K
  & layer temperatures, evaluated unclamped outside
  & Inside over the molecular layer \\
\end{longtable}

\section{Open items}

\begin{itemize}

\item \textbf{No $(v,J)$ resolution, and why it has not been added.} Carrying
  the ladder as transported species would mean $302$ further continuity
  equations, the $1833$-line radiative network as a level-coupling matrix
  instead of a summed emissivity, and a state-to-state collisional set that
  the code does not have and that none of its current sources supply. Against
  that cost, the place it would most obviously buy something --- the
  self-shielding --- is the place DB96 say the level distribution stops
  mattering, at columns above $\sim\!10^{18}$~cm$^{-2}$, three decades below
  where our layer sits. The trade is recorded here rather than settled.

\item \textbf{Self-shielding at our temperature is untested.} The
  insensitivity of \S\ref{sec:lw} was demonstrated among cold distributions.
  A calculation of $f_{\rm shield}$ on a $10^{3}$~K collisionally equilibrated
  population, at the columns of Table~\ref{tab:validity}, would settle whether
  the fit transfers; nothing in the literature consulted here does it.

\item \textbf{An internal inconsistency in the published DB96 fit.} DB96 cap
  the overlap-corrected equivalent width of the pumping lines at $W_{\max}
  \approx \ln(1110/912) \approx 0.2$ (their \S4.3, eqs.\ 29--30), i.e.\ at a
  band fraction of exactly $1$. Their closed-form integral, eq.~(39), has the
  asymptote $1.05(1 + 0.0117\,b_5)$ as $N_{\rm H_2} \rightarrow \infty$, which
  at our Doppler parameters is $1.0843$ ($b = 2.79$~km~s$^{-1}$) to $1.1050$
  ($4.48$~km~s$^{-1}$): $8.4$--$10.5\%$ above their own ceiling. Their remark
  below eq.~(39) is only that it ``corresponds to an equivalent width
  $W \approx \Delta\ln\nu \approx 0.2$ in the limit $N_2 \rightarrow
  \infty$''. The overshoot is a property of
  the published fit, not of our use of it, and is the same whether the band
  fraction is taken from eq.~(39) or from integrating $f_{\rm shield}$ on the
  grid. \texttt{fuv\_lw\_photon\_field} clamps it at $1$ so the transmission
  handed to the continuum absorbers of the same interval cannot go negative.

\item \textbf{The \texttt{mol\_ir\_bands} regression case does not test what
  its name says.} It sets \texttt{Base IR field: True} and
  \texttt{Molecular IR bands: True} but leaves \texttt{Oxygen chemistry} off,
  as its \texttt{EXHALE\_resolved.out} confirms, so there are no H$_2$O or CO
  carriers. Measured on
  \texttt{output/\allowbreak Cooling\_breakdown.txt}, columns 13 and 14 are exactly zero
  in all $504$ rows, while column 12, the H$_2$ line channel, is nonzero
  everywhere and peaks at $2.934\times10^{-7}$~erg~cm$^{-3}$~s$^{-1}$ at
  $r = 1.0025\,R_{\rm p}$, $T = 1316$~K. The zeros are structural rather than
  numerical: with $n({\rm H_2O}) = n({\rm CO}) = 0$ the two net-rate functions
  return $0$ on their own guard. The case gates the H$_2$ line channel and
  nothing else, and a case that exercises the H$_2$O and CO paths would have
  to add \texttt{Oxygen chemistry: True}.
  \emph{Note also} that the golden directory for this case carries only
  \texttt{Hydro\_ioniz.txt}, \texttt{Ion\_species.txt} and \texttt{run.log};
  \texttt{Cooling\_breakdown.txt} is not snapshotted, so the infrared columns
  are not under regression control at all.

\end{itemize}

\begin{thebibliography}{9}

\bibitem[Black \& Dalgarno (1977)]{BlackDalgarno1977}
J.~H.~Black \& A.~Dalgarno,
\emph{Models of interstellar clouds. I. The Zeta Ophiuchi cloud},
ApJS 34, 405 (1977).

\bibitem[Breshears \& Bird (1973)]{Breshears1973}
W.~D.~Breshears \& P.~F.~Bird,
\emph{Density gradient measurements of H$_2$ dissociation in shock waves},
14th Symposium (International) on Combustion, 211 (1973).

\bibitem[Cohen \& Westberg (1983)]{Cohen1983}
N.~Cohen \& K.~R.~Westberg,
\emph{Chemical kinetic data sheets for high-temperature chemical reactions},
J.~Phys.~Chem.~Ref.~Data 12, 531 (1983), p.~559.

\bibitem[Draine \& Bertoldi (1996)]{DB96}
B.~T.~Draine \& F.~Bertoldi,
\emph{Structure of stationary photodissociation fronts},
ApJ 468, 269 (1996).

\bibitem[Frelikh \& Murray-Clay (2026)]{Frelikh2026}
R.~Frelikh \& R.~A.~Murray-Clay,
\emph{Efficiency of hydrodynamic atmospheric escape in hot Jupiters and
super-Earths},
ApJ 996, 96 (2026).

\bibitem[Koskinen et al.\ (2022)]{Koskinen2022}
T.~T.~Koskinen et al.,
\emph{Mass loss by atmospheric escape from extremely close-in planets},
ApJ 929, 52 (2022).

\bibitem[Miller et al.\ (2013)]{Miller2013}
S.~Miller, T.~Stallard, J.~Tennyson \& H.~Melin,
\emph{Cooling by H$_3^+$ emission},
J.~Phys.~Chem.~A 117, 9770 (2013).

\bibitem[Roueff et al.\ (2019)]{Roueff2019}
E.~Roueff et al.,
\emph{The full infrared spectrum of molecular hydrogen},
A\&A 630, A58 (2019).

\bibitem[Yan et al.\ (1998)]{Yan1998}
M.~Yan, H.~R.~Sadeghpour \& A.~Dalgarno,
\emph{Photoionization cross sections of He and H$_2$},
ApJ 496, 1044 (1998).

\end{thebibliography}

\end{document}
